% explainer-source-hash: 758aa0e1fc4aa54b67ece6e4982367fb6dc2c725dd0193a8e15921f8ec2b6aa1
% explainer-source-commit: 208ef83ddf2d31892671bea435cbf0f1caa251fe
% explainer-generated: 2026-07-16
% Regenerate with /explain-all (see WIRING.md). Do not edit this stamp by hand:
% it is how the toolchain knows whether this document still matches the code.
\documentclass[11pt,a4paper]{article}
\input{preamble}

\title{\vspace{-1.5cm}\textbf{N8N Automations}\\[0.3em]
\large A Deep Dive into 59 Workflow Automations\\[0.2em]
\normalsize Exhaustive walkthrough, assuming no prior knowledge}
\author{}
\date{}

\begin{document}
\maketitle
\thispagestyle{empty}

\begin{keyidea}{What you are about to read}
This project is not a program in the usual sense. There is no source file you can run,
no test suite, no build. It is a collection of 59 \emph{automations}: business chores,
each one wired together as a flowchart inside a tool called n8n, then exported to a
file so it can be read, reviewed and re-imported.

This document explains what n8n is, what an exported workflow actually contains, what
each family of automations does and why, and where the collection is weak. Everything
here was checked against the exported files themselves, not taken on trust from the
project's own documentation.
\end{keyidea}

\tableofcontents
\newpage

\section{What this project is, and what it is not}

\subsection{The one-sentence version}

Fifty-nine business processes, each built as a flowchart that then runs itself, exported
as text files and documented honestly, faults included.

\subsection{What it is not}

It is worth being precise about this early, because the shape of the project is unusual
and it is easy to misjudge in both directions.

\begin{itemize}
  \item \textbf{It is not a codebase.} There is no application to compile. The only
    executable-ish artifacts are JavaScript snippets embedded inside individual
    workflow steps, and one small Python script used to draw pictures of the workflows
    for the website.
  \item \textbf{It is not a library or a product.} Nothing here is meant to be
    imported by another program. Each automation is a self-contained job.
  \item \textbf{It is not a live system.} These are \emph{exports}, snapshots taken out
    of a running n8n instance. Importing one into a fresh n8n gives you the shape of
    the automation but none of its account access, which has to be supplied again by
    hand.
  \item \textbf{It is not a set of toy examples.} Most of these ran, or were built to
    run, against real business accounts at a real agency: a real payment processor, a
    real customer database, real inboxes.
\end{itemize}

\begin{honest}{But the "runs itself" claim needs an immediate qualifier}
Of the 80 exported workflow files in this repository, exactly \textbf{9} carry
\code{active: true}. The other 71 are marked inactive. In n8n, an inactive workflow's
automatic triggers do not fire; it only runs when a human presses the button in the
editor.

That does not mean 71 automations are broken. An export is a snapshot, and a workflow
is routinely deactivated before export, or exported from a development copy. But it
does mean this repository is \textbf{not evidence that 59 automations are running in
production right now}, and nobody should describe it that way. It is evidence that 59
automations were designed and built. Section~\ref{sec:honest} returns to this.
\end{honest}

\section{n8n, explained from zero}

Nothing in this section assumes you have seen n8n, or any automation tool. If you
already know n8n well, skip to Section~\ref{sec:datamodel}.

\subsection{The problem automation tools solve}

Consider a real chore: \emph{when a sales deal is marked "won" in our customer database,
create a customer record in our payment processor and send them an invoice}.

You could do that by hand, every time, in about four minutes. You could also write a
program: talk to the database's API, talk to the payment processor's API, handle the
errors, host it somewhere, keep it alive. That is maybe a day of work and a permanent
maintenance burden, for four minutes of chore.

\begin{jargon}{API}
Short for Application Programming Interface. A door that a piece of software opens
deliberately so that \emph{other} software can talk to it, as opposed to the screens it
opens for humans. When this document says a workflow "calls Stripe's API", it means the
workflow sends Stripe a structured message over the internet asking it to do something,
and Stripe sends a structured answer back.
\end{jargon}

n8n sits between those two options. It is a \textbf{visual workflow automation tool}:
you drag boxes onto a canvas, wire them together with lines, and the tool runs the
resulting flowchart for you. It already knows how to talk to hundreds of services, so
"talk to the payment processor's API" becomes "drag the Stripe box onto the canvas".

\begin{keyidea}{The honest framing of what this skill is}
Dragging boxes sounds like it is not engineering. Building one of these badly is easy
and building one of these well is not, and the reason is that everything hard about
software is still present: state, failure, retries, idempotency, data shape mismatches,
rate limits, credentials, and the question of what happens on the run where the third
API call times out after the first two already took effect.

The visual editor removes the typing. It does not remove the thinking. Most of
Section~\ref{sec:honest} is a catalogue of places in this very repository where the
thinking was not done, and each one is a bug that would be equally possible in
handwritten code.
\end{keyidea}

\subsection{The five words you need}

\begin{jargon}{Node}
One box on the canvas: a single step. "Send a Gmail message" is a node. "Ask an AI
model for text" is a node. "If the amount is over 500, go left, otherwise go right" is a
node. A node takes data in, does one thing, and passes data out. Across this repository
there are \textbf{1{,}645} nodes.
\end{jargon}

\begin{jargon}{Connection}
A line between two nodes, saying "when this one finishes, run that one next, and hand it
your output". Connections are directional. The set of all connections is the flowchart.
\end{jargon}

\begin{jargon}{Trigger}
The special node that \emph{starts} a workflow. Every other node runs because something
before it ran; a trigger runs because something happened in the outside world. A
schedule trigger fires at 8am. A webhook trigger fires when another system sends this
workflow a message. A Telegram trigger fires when someone texts the bot. A manual
trigger fires when a human clicks "Execute" in the editor.
\end{jargon}

\begin{jargon}{Credential}
A stored login for an outside service, kept in n8n's own encrypted credential store and
referenced by nodes rather than typed into them. This distinction matters enormously for
this repository and Section~\ref{sec:security} is devoted to it.
\end{jargon}

\begin{jargon}{Expression}
A small piece of live text inside a node's settings that gets computed at run time
instead of being fixed. Written in double curly braces. If a node's "To" field contains
an expression referring to the email address from an earlier step, then each run sends
to a different person. A field without an expression is a hardcoded constant and sends
to the same person forever, which, as Section~\ref{sec:honest} shows, is a mistake this
repository makes more than once.
\end{jargon}

\subsection{The shape of a workflow}

Figure~\ref{fig:concept} shows the whole idea in one picture, using the deal-won example.

\begin{figure}[H]
\centering
\resizebox{\textwidth}{!}{
\begin{tikzpicture}[node distance=7mm]
  \node[extbox] (world) {Outside world\\\scriptsize a deal is marked "won"};
  \node[box, right=14mm of world, fill=warnred!12, draw=warnred] (trig) {\textbf{Trigger}\\\scriptsize Webhook};
  \node[box, right=10mm of trig] (fetch) {Node\\\scriptsize fetch the deal};
  \node[box, right=10mm of fetch] (cust) {Node\\\scriptsize create customer};
  \node[box, right=10mm of cust] (inv) {Node\\\scriptsize send invoice};

  \draw[flow] (world) -- (trig);
  \draw[flow] (trig) -- node[lbl, above] {data} (fetch);
  \draw[flow] (fetch) -- node[lbl, above] {data} (cust);
  \draw[flow] (cust) -- node[lbl, above] {data} (inv);

  \node[store, below=12mm of fetch] (cred) {credential\\store};
  \draw[dflow] (cred) -- (fetch);
  \draw[dflow] (cred) -- (cust);
  \draw[dflow] (cred) -- (inv);
  \node[lbl, below=1mm of cred] {logins live here, never in the flowchart};

  \node[lbl, above=8mm of cust, text width=6cm] {each box is a \textbf{node}; each arrow is a
    \textbf{connection}; the whole chain is a \textbf{workflow}};
\end{tikzpicture}}
\caption{The vocabulary in one picture. A trigger starts the run; nodes do one thing
each; connections carry data forward; credentials are referenced from a separate store
rather than embedded.}
\label{fig:concept}
\end{figure}

\subsection{Four more terms this repository leans on heavily}

\begin{jargon}{Sub-workflow}
A workflow called by another workflow, the way a function is called by other code. The
caller uses an \code{executeWorkflow} node; the callee starts with an
\code{executeWorkflowTrigger} node instead of a normal trigger. This repository uses the
pattern 17 times, and it is how the largest systems here are organized.
\end{jargon}

\begin{jargon}{AI Agent (and its tools)}
A node that hands a question to a large language model \emph{along with a list of things
the model is allowed to do}, and then lets the model decide which of them to do. Those
"things" are called tools, and each is another node plugged into the agent. An agent
with a calendar tool and an email tool can be asked "move my 3pm and tell Sam", and it
will work out that this means two tool calls in a particular order. There are
\textbf{109} agent nodes in this repository.
\end{jargon}

\begin{jargon}{LLM (Large Language Model)}
The kind of AI model behind ChatGPT and its competitors: it takes text and produces
plausible continuation text. Crucially for this document, an LLM's output is
\emph{probabilistic}, not guaranteed. Ask for JSON and you usually get JSON. "Usually"
is the entire source of a recurring class of bug in Section~\ref{sec:honest}.
\end{jargon}

\begin{jargon}{RAG (Retrieval-Augmented Generation)}
A technique for making an LLM answer from \emph{your} documents instead of its general
training. You chop the documents into chunks, convert each chunk into a list of numbers
capturing its meaning (an \emph{embedding}), and store them in a vector database. At
question time you embed the question, find the chunks whose numbers are nearest, and
paste those chunks into the prompt. The model then answers from real text you supplied
rather than from memory. Two systems here do this.
\end{jargon}

\section{The exported file: what is actually in it}
\label{sec:datamodel}

Every automation in this repository is, on disk, one or more \code{.json} files. This
section explains exactly what is inside one, because everything else in this document is
a reading of these files.

\begin{jargon}{JSON}
A plain-text format for structured data, made of nested lists and key/value pairs. It is
what most web services speak, and it is human-readable, which is what makes reviewing
these exports possible at all.
\end{jargon}

\subsection{The top level}

The export is one JSON object. The two keys that matter are \code{nodes} (a flat list of
every box on the canvas) and \code{connections} (a map describing every wire). The rest
is metadata: identifiers, timestamps, the workflow's display name, and the
\code{active} flag.

\begin{lstlisting}[language=,caption={The whole of a small export, with values that
identify a live endpoint replaced by placeholders. This is \file{youtube-transcript-fetcher},
the smallest automation here at 3 nodes.}]
{
  "name": "YT Transcript",
  "active": false,
  "nodes": [ ... 3 node objects ... ],
  "connections": {
    "Webhook":      { "main": [ [ { "node": "HTTP Request",       "index": 0 } ] ] },
    "HTTP Request": { "main": [ [ { "node": "Respond to Webhook", "index": 0 } ] ] }
  },
  "settings": { "executionOrder": "v1" },
  "pinData": {},
  "tags": []
}
\end{lstlisting}

\begin{keyidea}{Connections are keyed by node \emph{name}, not by ID}
Look closely at the \code{connections} map: the wires refer to nodes by their
human-readable display name (\code{"HTTP Request"}), even though every node also carries
a unique ID. The same is true of expressions inside nodes, which reach back to earlier
steps by writing \code{\$('Some Node Name')}.

This is the single most important structural fact about n8n for understanding
Section~\ref{sec:honest}, because it means \textbf{renaming a node in the editor can
silently break any expression that referred to it}. n8n fixes up the wires when you
rename, but an expression is just text inside a field, and text does not get fixed up.
This repository contains a workflow where exactly this happened, three times, in one
file.
\end{keyidea}

\subsection{One node, dissected}

\begin{lstlisting}[language=,caption={A single node object. This is the webhook trigger
from the same file; the real \code{path} and \code{webhookId} are the live URL of this
endpoint, so they are shown here as placeholders rather than reproduced.}]
{
  "name": "Webhook",                    <-- what connections and expressions call it
  "type": "n8n-nodes-base.webhook",     <-- which KIND of box this is
  "typeVersion": 2,                     <-- which generation of that box
  "position": [0, 0],                   <-- where it sits on the canvas (cosmetic)
  "id": "<uuid>",                       <-- unique, but NOT what wires use
  "webhookId": "<uuid>",
  "parameters": {                       <-- everything the user configured
    "httpMethod": "POST",
    "path": "<uuid>",
    "options": { "allowedOrigins": "*" }
  }
}
\end{lstlisting}

The \code{type} field is the interesting one. It is a namespaced string, and the
namespace tells you where the node came from:

\begin{center}
\begin{tabular}{@{}llr@{}}
\toprule
\textbf{Namespace} & \textbf{Meaning} & \textbf{Count} \\
\midrule
\code{n8n-nodes-base.*}              & Ships with n8n itself      & 1{,}298 \\
\code{@n8n/n8n-nodes-langchain.*}    & Official AI add-on pack    & 336 \\
\code{n8n-nodes-scrapegraphai.*}     & Third-party community pack & 4 \\
\code{n8n-nodes-fathom.*}            & Third-party community pack & 2 \\
\code{@apify/n8n-nodes-apify.*}      & Third-party community pack & 1 \\
\code{@devlikeapro/n8n-nodes-waha.*} & Third-party community pack & 1 \\
\code{@elevenlabs/n8n-nodes-elevenlabs.*} & Third-party community pack & 1 \\
\code{n8n-nodes-templated.*}         & Third-party community pack & 1 \\
\code{n8n-nodes-guuid-generator.*}   & Third-party community pack & 1 \\
\midrule
& & \textbf{1{,}645} \\
\bottomrule
\end{tabular}
\end{center}

\begin{honest}{Seven community packages are an undocumented portability constraint}
The bottom seven rows are not part of n8n. They are third-party packages that must be
installed into the instance before the workflow using them will even \emph{open}, let
alone run. Eleven nodes across seven separate packages, and each package is a separate
install and a separate maintenance risk: a community node abandoned by its author is a
workflow that stops opening after an n8n upgrade.

Anyone importing \file{fathom-native-node-integration},
\file{whatsapp-ai-sales-lead-system}, \file{uuid-generator-utility} or any of the four
others into a stock n8n gets an unrecognized-node error. Neither the repository root
README nor those folders' setup sections leads with that fact. It is a small thing that
will cost a first-time importer twenty confused minutes, and the fix is one line in each
affected \file{CREDENTIALS.md} or setup section.
\end{honest}

\subsection{How a wire is stored, and why the shape is weird}

\begin{lstlisting}[language=,caption={The nesting of the connections map, explained.}]
"connections": {
  "Loop Over Items": {          // the SOURCE node's name
    "main": [                   // which KIND of output (see below)
      [ {"node": "Switch2", "index": 0} ],      // output slot 0 -> these targets
      [ {"node": "Download File", "index": 0} ] // output slot 1 -> these targets
    ]
  }
}
\end{lstlisting}

Three levels of nesting, each meaning something:

\begin{enumerate}
  \item \textbf{The connection kind} (\code{"main"}). Ordinary data flow is
    \code{main}. AI nodes use other kinds: \code{ai\_languageModel},
    \code{ai\_tool}, \code{ai\_memory}, \code{ai\_embedding}. These are how a
    model, a tool or a memory gets attached to an agent.
  \item \textbf{The output slot} (the outer array). A node with one output has one
    slot. An \code{If} node has two: slot 0 is "true", slot 1 is "false". A
    \code{Switch} can have many. A loop node uses slot 0 for "done" and slot 1 for
    "next batch".
  \item \textbf{The target list} (the inner array). One output slot can feed several
    nodes at once, which is how the canvas fans out.
\end{enumerate}

\begin{keyidea}{The AI connections run \emph{backwards}, and this confuses everyone}
For an agent, the \code{ai\_tool} and \code{ai\_languageModel} connections are stored on
the \emph{tool}, pointing \emph{at} the agent. Reading the raw JSON, it looks like the
Gmail tool feeds the agent, in the same direction as ordinary data flow. It does not:
the agent calls the tool.

The arrow in the file describes attachment ("I am equipment belonging to that agent"),
not data flow. Anyone reading these exports for the first time, or writing a script to
analyse them, will get the direction wrong once. Figure~\ref{fig:assistant} draws it the
way it actually behaves rather than the way it is stored.
\end{keyidea}

\section{Repository anatomy}

\subsection{The layout}

\begin{figure}[H]
\centering
\begin{forest} dirtree
[n8n-automations/
  [README.md {\normalfont\itshape\quad the index, and the honest "what is missing" list}]
  [.gitignore]
  [crm-lead-to-stripe-invoice/ {\normalfont\itshape\quad one folder per automation, x59}
    [README.md {\normalfont\itshape\quad what/why/architecture/challenges}]
    [CREDENTIALS.md {\normalfont\itshape\quad which logins it needs, never their values}]
    [LICENSE {\normalfont\itshape\quad PolyForm Strict 1.0.0}]
    [workflow.json {\normalfont\itshape\quad the export itself}]
  ]
  [cold-email-outreach-sequence-system/ {\normalfont\itshape\quad some folders hold many}
    [README.md] [CREDENTIALS.md] [LICENSE]
    [Online\_Research.json]
    [NEW\_\_\_Generate\_All\_Emails.json {\normalfont\itshape\quad + 8 more}]
  ]
  [docs/
    [explainers/ {\normalfont\itshape\quad this document}]
    [screenshots/
      [render\_n8n.py {\normalfont\itshape\quad draws node graphs as PNGs}]
      [*.png {\normalfont\itshape\quad 6 rendered graphs}]
    ]
  ]
]
\end{forest}
\caption{The repository is one flat tier of automation folders plus a \file{docs/}
directory. There is no shared code, no top-level manifest, and no build.}
\label{fig:tree}
\end{figure}

\subsection{The four-file convention}

Every one of the 59 folders holds the same four things, and the consistency is real, not
aspirational. Verified by inspection:

\begin{center}
\begin{tabular}{@{}llp{7.4cm}@{}}
\toprule
\textbf{File} & \textbf{In all 59?} & \textbf{Role} \\
\midrule
\file{README.md}        & yes & The prose. Nine standard headings.\\
\file{CREDENTIALS.md}   & yes & The external accounts needed, by name and purpose.\\
\file{LICENSE}          & yes & PolyForm Strict 1.0.0, byte-identical in 58 of 59.\\
\file{*.json}           & yes & One or more exports. 80 files across 59 folders.\\
\bottomrule
\end{tabular}
\end{center}

The README headings are near-perfectly uniform across all 59 folders: \emph{What it is},
\emph{Why it exists}, \emph{Features}, \emph{Architecture}, \emph{Setup}, \emph{Usage},
\emph{Challenges}, \emph{What I learned}, \emph{What I'd do differently}. Two folders
lack a \emph{What it is} heading and one lacks \emph{Why it exists} and \emph{Features},
each having replaced it with something more specific such as \emph{Naming note (read
this first)}.

\begin{honest}{The 59th LICENSE differs, by one byte}
\file{uuid-generator-utility/LICENSE} is identical to the other 58 except that it lacks
a trailing newline. This is not a licensing problem and has no legal effect. It is worth
a single sentence only because it is a tell: the files were placed by hand or by a
script that did not normalize its output, rather than copied from one canonical source.
\end{honest}

\begin{jargon}{PolyForm Strict 1.0.0}
A source-available, not open-source, licence. Anyone may read the code; nobody may use
it, modify it or redistribute it without a separate agreement. The sensible choice for a
portfolio: the work can be examined by an interviewer but not lifted into a competitor's
product.
\end{jargon}

\section{The verified census}
\label{sec:census}

The project claims 59 automations. That claim is checkable and it was checked, by
enumerating the tracked files and parsing every export rather than by counting the
README's bullet points.

\begin{center}
\begin{tabular}{@{}lrl@{}}
\toprule
\textbf{Measure} & \textbf{Count} & \textbf{How it was established} \\
\midrule
Automation folders          & \textbf{59} & Tracked top-level directories, minus \file{docs/} \\
Exported workflow files     & 80          & Every tracked \file{*.json} \\
Total nodes                 & 1{,}645     & Sum over every export's \code{nodes} array \\
Distinct node types         & 93          & Unique \code{type} strings \\
Sticky notes (documentation)& 186         & \code{stickyNote} nodes \\
Functional nodes            & 1{,}459     & Total minus sticky notes \\
Workflows marked active     & 9 of 80     & \code{active: true} \\
Disabled nodes              & 45          & \code{disabled: true} \\
Workflows with pinned test data & 9       & Non-empty \code{pinData} \\
\bottomrule
\end{tabular}
\end{center}

\begin{keyidea}{The headline number is honest, with a caveat worth stating out loud}
\textbf{59 automations is correct}, if "automation" means "one documented folder solving
one business problem". It is the count of folders, and there are exactly 59.

The caveat is that 59 folders contain 80 workflow files, because six folders hold
multi-workflow systems (the cold email system alone is 10 files and 275 nodes). So "59"
undercounts the workflow files and overcounts the independent programs, depending on
which you meant. It is a defensible number, not an inflated one, and the alternative
counts are 80 files or 1{,}645 nodes.
\end{keyidea}

\begin{honest}{A different claim on the portfolio website is not about this project}
The website's project grid shows the metric "33 automation hooks" on a card. That card
belongs to \textbf{reelflow-studio}, a different project entirely, and the number has
nothing to do with this repository. This repository's card correctly reads "59
automations". Nothing here needs fixing; the number is noted only so it is not mistaken
for a claim this project failed to substantiate.
\end{honest}

\subsection{What the collection is built out of}

Node type frequency is a decent x-ray of a workflow collection's character.

\begin{figure}[H]
\centering
\begin{tikzpicture}
\begin{axis}[
  width=13.5cm, height=7.4cm,
  xbar, xmin=0,
  bar width=8pt,
  y=10pt,
  enlarge y limits=0.06,
  xlabel={\small number of nodes across all 80 exports},
  symbolic y coords={agent,gmail,wait,openAi,merge,code,if,lmChatGoogleGemini,googleSheets,httpRequest,set,stickyNote},
  ytick=data,
  nodes near coords, nodes near coords style={font=\tiny},
  tick label style={font=\scriptsize},
  label style={font=\small},
  every axis plot/.append style={fill=accent!55, draw=inkblue},
  axis lines*=left,
  xmajorgrids, grid style={linegrey},
]
\addplot coordinates {
  (186,stickyNote) (136,set) (128,httpRequest) (124,googleSheets)
  (109,agent) (82,lmChatGoogleGemini) (80,if) (71,code)
  (56,merge) (37,openAi) (33,wait) (32,gmail)
};
\end{axis}
\end{tikzpicture}
\caption{The twelve most common node types. Read it as a portrait: this is a collection
about \emph{moving data between web services and an LLM}, glued with spreadsheets.}
\label{fig:nodefreq}
\end{figure}

Four things fall straight out of that chart:

\begin{itemize}
  \item \textbf{Sticky notes are the single most common node} (186). A sticky note does
    nothing at run time; it is a comment on the canvas. Roughly one node in nine here
    exists purely to explain the others. That is a genuinely good sign about how these
    were built, and it is unusual.
  \item \textbf{Google Sheets is the database} (124 nodes, plus 5 Sheets triggers).
    Section~\ref{sec:patterns} argues this is both the collection's cleverest habit and
    its biggest architectural weakness.
  \item \textbf{\code{httpRequest} appears 128 times}, meaning that for a large fraction
    of the services touched, no purpose-built n8n node existed and the author wrote the
    raw API call by hand. This is the "when the visual tool runs out, drop to the API"
    move, and it is used constantly.
  \item \textbf{AI is not a garnish here.} 109 agent nodes plus 82 Gemini models plus
    37 OpenAI nodes plus 17 OpenAI chat models. A fifth of all nodes are AI-related.
\end{itemize}

\subsection{How these automations start}

\begin{center}
\begin{tabular}{@{}lrp{8.1cm}@{}}
\toprule
\textbf{Trigger type} & \textbf{n} & \textbf{Meaning} \\
\midrule
\code{manualTrigger}          & 28 & A human clicks Execute \\
\code{executeWorkflowTrigger} & 16 & Called by another workflow \\
\code{scheduleTrigger}        & 14 & Fires on a clock \\
\code{webhook}                & 14 & Another system posts to a URL \\
\code{formTrigger}            & 11 & Someone submits an n8n-hosted form \\
\code{chatTrigger}            & 7  & Someone types into a chat box \\
\code{telegramTrigger}        & 5  & Someone messages a Telegram bot \\
\code{googleSheetsTrigger}    & 5  & A spreadsheet row changed \\
\code{googleDriveTrigger}     & 3  & A file appeared or changed \\
\code{errorTrigger}           & 2  & Another workflow failed \\
\code{fathomTrigger}          & 2  & A meeting recording finished \\
\code{calTrigger}             & 1  & A booking was made \\
\code{wahaTrigger}            & 1  & A WhatsApp message arrived \\
\bottomrule
\end{tabular}
\end{center}

\begin{honest}{28 manual triggers is the most revealing number in this document}
\code{manualTrigger} is the most common trigger in the collection. A manual trigger is
what n8n gives you while you are \emph{building}: it exists so you can press Execute and
watch data flow. A finished, deployed automation is normally started by a schedule, a
webhook or an event, not by a person opening the editor.

Combined with only 9 of 80 workflows marked active, the fair reading is that a
substantial share of this collection is at \textbf{"built and demonstrated"} maturity
rather than \textbf{"deployed and running unattended"} maturity. Several of the folder
READMEs say as much. It should be said plainly and first, rather than discovered by an
interviewer who thought to count.

This is not a criticism of the work. It is a criticism of any description of the work
that implies 59 live production systems.
\end{honest}

\section{A taxonomy of the 59}

Fifty-nine is too many to walk through one at a time, and they are not 59 unrelated
things. They fall into seven families, and once you know the family you can predict most
of the architecture.

\begin{figure}[H]
\centering
\resizebox{\textwidth}{!}{
\begin{tikzpicture}[node distance=5mm]
  \node[box, fill=inkblue!12, minimum width=3.2cm] (root) {\textbf{59 automations}};

  \node[box, below left=12mm and 34mm of root, text width=3.0cm] (a) {\textbf{Lead generation\\and outreach}\\\scriptsize 11 folders};
  \node[box, right=6mm of a, text width=3.0cm] (b) {\textbf{Content\\generation}\\\scriptsize 14 folders};
  \node[box, right=6mm of b, text width=3.0cm] (c) {\textbf{Video\\production}\\\scriptsize 7 folders};
  \node[box, right=6mm of c, text width=3.0cm] (d) {\textbf{SEO}\\\scriptsize 5 folders};

  \node[box, below=16mm of a, text width=3.0cm] (e) {\textbf{Assistants\\and chatbots}\\\scriptsize 9 folders};
  \node[box, right=6mm of e, text width=3.0cm] (f) {\textbf{Business ops}\\\scriptsize 6 folders};
  \node[box, right=6mm of f, text width=3.0cm] (g) {\textbf{Utilities}\\\scriptsize 7 folders};
  \node[extbox, right=6mm of g, text width=3.0cm] (h) {\scriptsize dominant shape:\\\scriptsize trigger $\rightarrow$ fetch $\rightarrow$\\\scriptsize LLM $\rightarrow$ deliver\\\scriptsize $\rightarrow$ log to Sheets};

  \foreach \x in {a,b,c,d} { \draw[flow] (root) -- (\x.north); }
  \foreach \x in {e,f,g} { \draw[flow] (root) to[out=270,in=90] (\x.north); }
\end{tikzpicture}}
\caption{The seven families. Counts are of folders and sum to 59.}
\label{fig:taxonomy}
\end{figure}

\subsection{Family 1: Lead generation and outreach (11)}

Find businesses, get contact details, write to them, chase them.

\begin{flushleft}
\file{google-maps-to-email-scraper}, \file{linkedin-extractor-from-google-maps},
\file{instagram-link-finder-outreach}, \file{cold-email-outreach-sequence-system},
\file{deep-multiline-icebreaker-system}, \file{gmail-outreach-follow-ups},
\file{subscriber-email-sequence}, \file{newsletter-automation},
\file{whatsapp-ai-sales-lead-system}, \file{upwork-ai-agent-pipeline},
\file{instagram-viral-reel-scraper}.
\end{flushleft}

The shared shape: a scraping API supplies raw businesses, a cleaning step normalizes
them, an LLM writes something personal per lead, a sending platform delivers it, and a
spreadsheet records who has been contacted. This family contains the largest system here
and Section~\ref{sec:coldemail} walks it.

\subsection{Family 2: Content generation (14)}

\begin{flushleft}
\file{social-media-content-generator} (190 nodes, the biggest single file),
\file{ai-social-media-amplifier}, \file{social-media-poster},
\file{social-media-trend-monitor}, \file{instagram-posting-workflow},
\file{linkedin-publish-utility}, \file{real-estate-market-update-poster},
\file{fb-ad-creative-generator}, \file{text-to-image-generator},
\file{presentation-generator-tool}, \file{ai-proposal-generator-system},
\file{daily-health-tip-bot}, \file{ai-motivational-quotes-bot},
\file{ai-workout-generator}.
\end{flushleft}

Shape: gather a fact or a prompt, ask an LLM to write copy, optionally ask an image
model for a picture, publish or return it.

\subsection{Family 3: Video production (7)}

\begin{flushleft}
\file{auto-video-generator}, \file{faceless-shorts-system},
\file{heygen-ai-video-production}, \file{submagic-video-clip-generator},
\file{instagram-reel-analyzer}, \file{instagram-reel-scraper-transcriber},
\file{youtube-agent-starter-template}.
\end{flushleft}

Shape: script $\rightarrow$ voice $\rightarrow$ render $\rightarrow$ poll $\rightarrow$
publish. This family is defined by the polling problem: video rendering takes minutes,
so these workflows must ask "is it done yet?" on a loop, which is why 33 \code{wait}
nodes exist in this repository.

\subsection{Family 4: SEO (5)}

\begin{flushleft}
\file{seo-ai-agent}, \file{seo-keyword-research-tool},
\file{cluster-seo-post-generator}, \file{trend-seo-post-generator},
\file{competitor-content-gap-analyzer}.
\end{flushleft}

Shape: a keyword data API supplies volumes, an LLM clusters and writes, WordPress
receives a draft. Notably these publish as \emph{drafts}, never live, which is the right
call and is stated as deliberate.

\subsection{Family 5: Assistants and chatbots (9)}

\begin{flushleft}
\file{ultimate-personal-assistant-system}, \file{rag-ingestion-pipeline},
\file{tutorial-rag-chat-agent}, \file{content-business-chat-agent},
\file{telegram-joke-agent}, \file{telegram-translator-bot},
\file{telegram-checklist-generator}, \file{telegram-to-vapi-trigger},
\file{deep-research-report-generator}.
\end{flushleft}

Shape: a chat or messaging trigger, an agent node, a memory node, and a set of tools.
Sections~\ref{sec:assistant} and~\ref{sec:rag} walk the two most interesting.

\subsection{Family 6: Business operations (6)}

\begin{flushleft}
\file{crm-lead-to-stripe-invoice}, \file{calcom-trigger-workflow},
\file{fathom-native-node-integration}, \file{expense-tracker-to-sheets},
\file{daily-email-digest-ai-powered}, \file{birthday-greetings-automation}.
\end{flushleft}

Shape: a real business event happens, and something is recorded, billed or announced.
This is the family where a bug costs actual money, and Section~\ref{sec:invoice} walks
the one that bills people.

\subsection{Family 7: Utilities (7)}

\begin{flushleft}
\file{uuid-generator-utility}, \file{website-screenshot-utility},
\file{youtube-transcript-fetcher}, \file{handwriting-to-text-ocr},
\file{booking-time-utility-set}, \file{vapi-call-trigger-utility},
\file{hacker-news-daily-digest}.
\end{flushleft}

Small, single-purpose, mostly webhook-in / answer-out. These are the building blocks
other workflows call, and the smallest is 3 nodes.

\section{Walkthrough 1: CRM lead to Stripe invoice}
\label{sec:invoice}

This is the automation to understand first. It is only 6 nodes, it is the example from
Section~2, and it touches real money, which makes its flaws the most instructive in the
repository.

\subsection{What it does}

When a sales task in ClickUp (a project management tool used here as a CRM) is moved to
the closed-won status, a ClickUp automation posts to this workflow's webhook. The
workflow then fetches the full task, creates a Stripe customer, and issues a finalized
invoice.

\begin{figure}[H]
\centering
\resizebox{\textwidth}{!}{
\begin{tikzpicture}[node distance=6mm]
  \node[extbox, text width=2.3cm] (cu) {ClickUp\\\scriptsize deal marked won};
  \node[box, right=9mm of cu, fill=warnred!12, draw=warnred, text width=1.9cm] (wh) {Webhook\\\scriptsize trigger};
  \node[box, right=9mm of wh, text width=2.1cm] (get) {ClickUp\\\scriptsize get task by id};
  \node[box, right=9mm of get, text width=2.2cm] (st) {Stripe\\\scriptsize create customer};
  \node[box, right=9mm of st, text width=2.2cm] (ii) {HTTP\\\scriptsize invoice item};
  \node[box, right=9mm of ii, text width=2.1cm] (ci) {HTTP\\\scriptsize create invoice};
  \node[box, right=9mm of ci, text width=2.1cm] (fi) {HTTP\\\scriptsize finalize};

  \draw[flow] (cu) -- (wh);
  \draw[flow] (wh) -- (get);
  \draw[flow] (get) -- (st);
  \draw[flow] (st) -- (ii);
  \draw[flow] (ii) -- (ci);
  \draw[flow] (ci) -- (fi);

  \node[lbl, below=7mm of get, text width=3.2cm, draw=warnred, fill=warnred!8] (b1)
    {\textbf{email} read as\\ \code{custom\_fields[19]}};
  \draw[dflow, draw=warnred] (b1) -- (st);
  \node[lbl, below=7mm of ii, text width=3.4cm, draw=warnred, fill=warnred!8] (b2)
    {\textbf{amount} read as\\ \code{custom\_fields[5] * 100}\\ description hardcoded};
  \draw[dflow, draw=warnred] (b2) -- (ii);

  \node[lbl, above=6mm of ci, text width=5.2cm] {no branch, no retry, no idempotency check:\\ any repeat delivery bills the customer twice};
\end{tikzpicture}}
\caption{Six nodes, one straight line, no error path. The two red annotations are
verified readings of the actual node parameters, not hypotheticals.}
\label{fig:invoice}
\end{figure}

\subsection{Why three raw HTTP calls sit next to a native Stripe node}

The workflow uses n8n's built-in Stripe node to create the customer, and then three
hand-written HTTP calls to do the invoicing. That looks inconsistent, and it is worth
explaining because it is the single most common shape in this repository.

n8n's Stripe node supports creating customers. It does not support the three-step
invoice dance Stripe requires (create an invoice \emph{item}, create the \emph{invoice},
then \emph{finalize} it to actually send it). So the author used the native node where
it fit, and dropped to raw API calls where it did not, reusing the same stored Stripe
credential for both via n8n's \code{predefinedCredentialType} option so the API key
still never appears in the workflow.

\begin{keyidea}{This is the defining move of the whole collection}
128 \code{httpRequest} nodes exist for this reason. The visual tool covers the common
80\%; the remaining 20\% is done by reading the vendor's API docs and constructing the
call by hand. Someone who can only drag boxes stops at the edge of what the boxes do.
Being able to say why each of those 128 exists is what makes this a portfolio project
rather than a folder of exports.
\end{keyidea}

\subsection{The flaws, verified}

Every one of the following was confirmed by reading the node parameters directly, not
inferred from the README.

\begin{honest}{Custom fields are read by array position}
The Stripe node's email is the expression \code{custom\_fields[19].value}. The amount is
\code{custom\_fields[5].value * 100}.

ClickUp returns custom fields as an ordered array. Position 19 is the email
\emph{today}. Add a field to that ClickUp list, reorder the columns, or delete one, and
19 silently becomes something else. There is no name lookup, no assertion, no guard. The
workflow does not fail; it invoices the wrong data to the wrong address, quietly, and
the first person to notice is the customer.

Fixing it is a small code node that finds the field by its name or ID. The fact that the
README identifies this flaw itself is genuinely to the author's credit; it is still
present in the file.
\end{honest}

\begin{honest}{No idempotency, on a node that charges people}
The chain has no check for an existing Stripe customer with that email, and no check for
an existing invoice for that deal. Webhooks are retried by their senders as a matter of
routine when a response is slow. Two deliveries of the same event produce two customers
and two invoices.

This is the highest-severity issue in the repository. Everything else here produces a
bad tweet or a wasted API call; this one double-bills a real person. The standard fix is
an idempotency key, which Stripe explicitly supports on exactly these endpoints.
\end{honest}

\begin{honest}{The invoice line always says the same thing}
The invoice item's description is the literal string \code{"1st Deposit"}, not an
expression. Every invoice this workflow has ever created carries that line, whatever the
deal actually was.
\end{honest}

\section{Walkthrough 2: the Ultimate Personal Assistant}
\label{sec:assistant}

Five workflow files, 50 nodes. This is the collection's cleanest architecture and the
best demonstration of the sub-workflow-as-tool pattern.

\subsection{What it does}

You text a Telegram bot, by typing or by voice. An orchestrating agent works out what
you want and delegates to one of four specialist agents: email, calendar, contacts, or
content creation. The answer comes back as a Telegram reply.

\begin{figure}[H]
\centering
\resizebox{\textwidth}{!}{
\begin{tikzpicture}[x=1cm, y=1cm]
  % --- input column ---
  \node[box, fill=warnred!12, draw=warnred, text width=2.0cm] (tg) at (0,0) {Telegram\\Trigger};
  \node[dec] (sw) at (3.1,0) {text or\\voice?};
  \node[box, text width=2.0cm] (set) at (6.4,1.3) {Set 'Text'};
  \node[box, text width=2.0cm] (dl)  at (6.4,-1.3) {Download\\File};
  \node[box, text width=2.0cm] (tr)  at (9.2,-1.3) {Transcribe\\\scriptsize Whisper};

  % --- the agent ---
  \node[box, fill=accent!18, minimum height=15mm, text width=2.5cm] (ag) at (12.4,0)
    {\textbf{Ultimate}\\\textbf{Assistant}\\\scriptsize agent};

  \draw[flow] (tg) -- (sw);
  \draw[flow] (sw.north) |- node[lbl,pos=0.72,above]{text} (set.west);
  \draw[flow] (sw.south) |- node[lbl,pos=0.72,below]{voice} (dl.west);
  \draw[flow] (dl) -- (tr);
  \draw[flow] (set.east) to[out=0,in=160] (ag.160);
  \draw[flow] (tr.east)  to[out=0,in=200] (ag.200);

  % --- equipment, below the agent ---
  \node[box, text width=2.1cm] (llm) at (11.1,-3.4) {OpenAI\\Chat Model};
  \node[box, text width=2.1cm] (mem) at (13.9,-3.4) {Window Buffer\\Memory};
  \draw[dflow] (llm.north) -- (ag.250);
  \draw[dflow] (mem.north) -- (ag.290);
  \node[lbl] at (12.5,-4.4) {equipment attached to the agent};

  % --- tools, right column ---
  \node[box, text width=2.4cm] (t1) at (17.0,2.7)  {Email Agent};
  \node[box, text width=2.4cm] (t2) at (17.0,1.5)  {Calendar Agent};
  \node[box, text width=2.4cm] (t3) at (17.0,0.3)  {Contact Agent};
  \node[box, text width=2.4cm] (t4) at (17.0,-0.9) {Content Creator};
  \node[extbox, text width=2.4cm] (t5) at (17.0,-2.1) {Tavily (web)};
  \node[extbox, text width=2.4cm] (t6) at (17.0,-3.3) {Calculator};

  \foreach \t in {t1,t2,t3,t4,t5,t6} { \draw[flow, draw=accent] (ag.east) -- (\t.west); }

  \begin{scope}[on background layer]
    \node[fit=(t1)(t4), draw=accent, dashed, rounded corners, inner sep=4pt] (grp) {};
  \end{scope}
  \node[lbl, text width=3.1cm, align=left, anchor=west] at (19.0,0.9)
    {each of these four is a\\\textbf{separate workflow file},\\called as a tool};

  % --- reply ---
  \node[box, text width=2.0cm] (resp) at (12.4,4.4) {Telegram\\Response};
  \draw[flow] (ag.north) to[out=100,in=270] (resp.south);
  \node[lbl, anchor=west] at (13.6,4.4) {the agent's answer goes back to the chat};
\end{tikzpicture}}
\caption{The orchestrator. Solid blue arrows from the agent are tool \emph{calls} (the
agent decides). Dashed arrows are equipment attachments. Note that the file stores the
tool arrows pointing the other way; this diagram shows behaviour, not storage.}
\label{fig:assistant}
\end{figure}

\subsection{Why this design is good}

\begin{keyidea}{Sub-workflow as tool is the automation equivalent of a function call}
Each specialist is its own file with its own \code{executeWorkflowTrigger}. That buys
exactly what functions buy in code:

\begin{itemize}
  \item \textbf{Testability.} The Calendar Agent can be executed on its own with pinned
    input, without going near Telegram.
  \item \textbf{Isolation.} Changing how contacts are looked up cannot break email.
  \item \textbf{Prompt size.} A single agent holding calendar, email, contacts and
    content instructions would carry an enormous system prompt and choose badly. Four
    focused agents each hold a small one.
  \item \textbf{Reuse.} The Email Agent is callable by anything else.
\end{itemize}

The alternative, which several other workflows in this repository chose, is one giant
canvas. Compare the 190-node single file in \file{social-media-content-generator}.
\end{keyidea}

The voice handling deserves a note: the Switch on the trigger's output distinguishes a
text message from a voice note, and the voice branch downloads the audio and pushes it
through a transcription node before rejoining the same agent. Both paths converge, so
the agent never knows or cares which one it was. That is the correct place to put that
seam.

\begin{honest}{The orchestrator is 15 nodes; the intelligence is in prompts you cannot see}
Structurally this file is nearly empty: a trigger, a switch, two branches, an agent, six
tools, a reply. Everything that decides \emph{behaviour} lives in the agent's system
prompt and in the four sub-agents' prompts.

That is not a flaw in the design, but it is a limit on what this document can tell you.
Reading the graph proves the wiring is sound. It cannot prove the assistant is any good,
and no amount of static reading will. That would need execution logs, which are not in
the repository.
\end{honest}

\section{Walkthrough 3: the RAG ingestion pipeline}
\label{sec:rag}

47 nodes. The most technically dense single automation here, and the one that best
rewards a diagram.

\subsection{What it does}

Watch a Google Drive folder. When a document appears, work out its type, extract its
text, chop it up, turn the chunks into embeddings, and store them in a Supabase vector
table. Separately, a chat agent answers questions using that table as its source.

\begin{figure}[H]
\centering
\resizebox{\textwidth}{!}{
\begin{tikzpicture}[x=1cm, y=1cm]
  % ---------- ingestion lobe ----------
  \node[box, fill=warnred!12, draw=warnred, text width=2.1cm] (fc) at (0,0.9) {File Created\\\scriptsize Drive trigger};
  \node[box, dashed, text width=2.1cm] (fu) at (0,-0.9) {File Updated\\\scriptsize \textbf{DISABLED}};
  \node[box, text width=1.8cm] (sid)  at (3.2,0.9) {Set\\File ID};
  \node[box, text width=2.0cm] (loop) at (6.2,0.9) {Loop Over\\Items};
  \node[box, text width=1.8cm] (dlf)  at (9.2,0.9) {Download\\File};
  \node[dec] (swi) at (11.9,0.9) {Switch\\type?};

  \node[box, text width=2.0cm] (pdf) at (15.2,3.3)  {Extract PDF};
  \node[box, text width=2.0cm] (txt) at (15.2,1.9)  {Extract Text};
  \node[box, text width=2.0cm] (xls) at (15.2,0.4)  {Extract Excel};
  \node[box, text width=2.0cm] (gd)  at (15.2,-1.2) {Convert to\\Google Doc};

  \node[box, text width=1.9cm] (agg) at (18.2,0.4) {Aggregate};
  \node[box, text width=1.9cm] (sum) at (21.0,0.4) {Summarize};

  \node[box, fill=accent!18, text width=2.2cm, minimum height=13mm] (vs) at (24.2,2.6)
    {Insert into\\Supabase\\Vectorstore};
  \node[store, text width=1.7cm] (sb) at (27.6,2.6) {Supabase\\vectors};

  \draw[flow] (fc) -- (sid);
  \draw[flow, dashed] (fu.east) -| (sid.south);
  \draw[flow] (sid) -- (loop);
  \draw[flow] (loop) -- node[lbl,above]{batch} (dlf);
  \draw[flow] (dlf) -- (swi);
  \draw[flow] (swi.east) to[out=0,in=180] (pdf.west);
  \draw[flow] (swi.east) to[out=0,in=180] (txt.west);
  \draw[flow] (swi.east) to[out=0,in=180] (xls.west);
  \draw[flow] (swi.east) to[out=0,in=180] (gd.west);
  \draw[flow] (xls) -- (agg);
  \draw[flow] (agg) -- (sum);
  \draw[flow] (pdf.east) to[out=0,in=180] (vs.west);
  \draw[flow] (txt.east) to[out=0,in=190] (vs.210);
  \draw[flow] (sum.east) to[out=0,in=250] (vs.250);
  \draw[flow] (vs) -- (sb);

  % equipment under the vector store
  \node[box, text width=2.0cm] (spl) at (23.0,0.2) {Text\\Splitter};
  \node[box, text width=2.0cm] (emb) at (25.8,0.2) {Embeddings\\OpenAI};
  \draw[dflow] (spl.north) -- (vs.290);
  \draw[dflow] (emb.north) -- (vs.330);

  % ---------- query lobe ----------
  \node[box, fill=warnred!12, draw=warnred, text width=2.2cm] (chat) at (14.0,-4.6) {Chat message\\received};
  \node[box, fill=accent!18, text width=2.0cm, minimum height=12mm] (rag) at (18.0,-4.6) {New RAG\\Agent};
  \node[box, text width=2.0cm] (pg) at (18.0,-6.4) {Postgres\\Memory};
  \node[box, text width=2.0cm] (vst) at (22.0,-4.6) {Vector Store\\Tool};
  \draw[flow] (chat) -- (rag);
  \draw[dflow] (pg.north) -- (rag.south);
  \draw[flow, draw=accent] (rag) -- (vst);
  \draw[flow] (vst.east) to[out=0,in=270] (sb.south);

  \begin{scope}[on background layer]
    \node[fit=(fc)(fu)(gd)(sb)(spl)(emb), draw=linegrey, rounded corners,
          inner sep=7pt, fill=softgrey!50] (ing) {};
    \node[fit=(chat)(vst)(pg), draw=linegrey, rounded corners,
          inner sep=7pt, fill=softgrey!50] (qry) {};
  \end{scope}
  \node[lbl, anchor=west, font=\small\itshape] at (ing.north west) {\textbf{INGESTION: get documents in}};
  \node[lbl, anchor=west, font=\small\itshape] at (qry.south west) {\textbf{QUERY: answer from them}};
\end{tikzpicture}}
\caption{Two independent lobes sharing one vector table. The dashed "File Updated" path
is a disabled trigger feeding a duplicate copy of the entire ingestion branch.}
\label{fig:rag}
\end{figure}

\subsection{The interesting parts}

\textbf{The type switch.} Documents arrive as PDFs, plain text, spreadsheets and Google
Docs, and each needs different extraction. The Switch node routes six ways: three go to
dedicated extractors, and the other three all converge on "convert to a Google Doc
first, then handle it". That is a nice reduction: rather than write four more extractors,
push the odd formats through a format you already handle.

\textbf{The spreadsheet detour.} Excel files do not go straight to the vector store; they
go through Aggregate and then Summarize first. This is correct and shows real
understanding of RAG. Embedding a spreadsheet row by row produces chunks like
\code{"Q3 \textbar{} 41 \textbar{} 12"}, which carry no retrievable meaning. Summarizing
into prose first gives the embedding something to work with.

\textbf{The delete-before-insert.} On the update path, \code{Delete Old Doc Rows} runs
before re-inserting. Without it, editing a document would leave the old chunks in the
table alongside the new ones, and the agent would answer from both, confidently citing
a version that no longer exists. Someone thought about this.

\begin{honest}{The entire ingestion branch exists twice, and the second copy is switched off}
The graph contains two complete, parallel ingestion pipelines. The \code{File Created}
trigger drives one (Set File ID1, Loop Over Items, Switch, Extract PDF Text1, Extract
from Text File1, Extract from Excel1, Aggregate, Summarize, Insert into Supabase
Vectorstore). The \code{File Updated} trigger drives a second, near-identical copy with
different node names (Set File ID, Loop Over Items1, Switch2, Extract PDF Text, ...,
Insert into Supabase Vectorstore1). The \code{File Updated} trigger is
\textbf{disabled}.

So: roughly half of this 47-node workflow is a duplicate of the other half, and it is
currently dead. Every future change to the extraction logic has to be made twice, in two
places, with confusingly similar node names, and a mistake in the disabled copy will not
show up until someone re-enables it months later.

The clean fix exists and is not exotic: both triggers feed a single shared branch, with
the delete step gated behind a check for which trigger fired. This is copy-paste
architecture, and it is the most expensive kind of debt in a visual tool, because the
tool makes copying a canvas trivially easy and offers nothing that would make you factor
it out.
\end{honest}

\section{Walkthrough 4: the cold email outreach system}
\label{sec:coldemail}

Ten files, 275 nodes. The largest system here by a wide margin, and the one that best
shows both the ambition and the structural problems of building at this scale in a
visual tool.

\subsection{What it does}

Take a raw lead (name, company, LinkedIn URL). Research them. Write a six-touch email
sequence personalized to what the research found. Push the finished sequence into
Instantly.ai, a cold-email sending platform.

\begin{figure}[H]
\centering
\resizebox{\textwidth}{!}{
\begin{tikzpicture}[node distance=5mm]
  \node[box, fill=warnred!12, draw=warnred, text width=2.2cm] (res) {\textbf{Online Research}\\\scriptsize manual trigger};
  \node[extbox, right=8mm of res, text width=2.2cm] (tav) {Tavily API\\\scriptsize x4 lookups};
  \draw[flow] (res) -- (tav);
  \node[store, right=8mm of tav, text width=2.0cm] (sheet) {Google\\Sheet};
  \draw[flow] (tav) -- node[lbl, above, align=center]{4 fields} (sheet);

  \node[lbl, below=2mm of tav, text width=4.6cm, draw=warnred, fill=warnred!8]
    {\textbf{no wire to the orchestrator.}\\the handoff is the sheet, and it is \emph{inferred} from matching field names};

  \node[box, fill=warnred!12, draw=warnred, below=24mm of res, text width=2.4cm] (orch)
    {\textbf{Generate All Emails}\\\scriptsize manual trigger, loops leads};
  \draw[dflow, draw=warnred] (sheet) to[out=270,in=0] (orch);

  \node[box, right=30mm of orch, text width=2.3cm] (s1) {First Email};
  \node[box, above=2mm of s1, text width=2.3cm] (s0) {Subject Line};
  \node[box, below=2mm of s1, text width=2.3cm] (s2) {Follow Up 1};
  \node[box, below=2mm of s2, text width=2.3cm] (s3) {Follow Up 2};
  \node[box, below=2mm of s3, text width=2.3cm] (s4) {Follow Up 3};
  \node[box, below=2mm of s4, text width=2.3cm] (s5) {Follow Up 4};
  \node[box, below=2mm of s5, text width=2.3cm] (s6) {Follow Up 5};
  \node[box, below=2mm of s6, text width=2.3cm] (s7) {Add to Instantly};

  \foreach \s in {s0,s1,s2,s3,s4,s5,s6,s7} { \draw[flow] (orch) -- (\s.west); }

  \node[extbox, right=9mm of s3, text width=2.2cm, minimum height=20mm] (inst) {Instantly.ai\\\scriptsize 5 campaigns\\\scriptsize routed by type};
  \draw[flow] (s7) -| (inst);

  \node[lbl, above left=1mm and 1mm of s1.west, text width=2.6cm, align=right]
    {8 sub-workflow\\calls per lead};

  \begin{scope}[on background layer]
    \node[fit=(s0)(s6), draw=accent, dashed, rounded corners, inner sep=5pt] (grp) {};
  \end{scope}
  \node[lbl, right=3mm of grp.north east, anchor=north west, text width=5.2cm, align=left]
    {each of these 6 runs \textbf{5 parallel LLM attempts} internally, then picks one};
\end{tikzpicture}}
\caption{The cold email system. Up to 30 Gemini calls per lead before a single sequence
is finished.}
\label{fig:coldemail}
\end{figure}

\subsection{The five-attempt pattern}

Each of the six copywriting workflows generates the same email five times in parallel
(\code{generate\_v1} through \code{generate\_v5}), each with its own model node, then
uses a switch/if gate to pick the first non-empty result.

\begin{keyidea}{Why generate the same thing five times?}
Because an LLM call can come back empty, refuse, or return unusable formatting, and it
does so unpredictably. Running one attempt and retrying on failure costs you the
latency of a second round trip. Running five in parallel and taking the first good one
costs you five times the money and none of the time.

It is a legitimate engineering trade, and it is the right one when a human is waiting.
It is the wrong one when nobody is waiting, which is this case: this is a batch job
over a lead list. The five-way fan-out buys latency that nobody is spending, at 5x the
token cost. A single call with a retry-on-empty would be materially cheaper for the
same result.
\end{keyidea}

\begin{honest}{Six copies of the same logic, kept in sync by hand}
The five-attempt block and its output-cleaning code node are duplicated across all six
copywriting workflows. There is no shared sub-workflow for the common part, even though
this system \emph{already uses} sub-workflows and therefore had the mechanism sitting
right there.

A fix to the cleaning regex has to be made six times. The consequences of that are
already visible in the repository: \file{NEW\_\_\_Generate\_All\_Emails.json} still
contains \textbf{nine orphaned \code{executeWorkflow} nodes} from an earlier version,
disconnected from everything, left behind when the workflow was restructured. Verified:
they have no incoming and no outgoing connections.
\end{honest}

\begin{honest}{The research step is not actually connected to anything}
\file{Online\_Research.json} has a manual trigger and is called by no
\code{executeWorkflow} node in any of the other nine files. Its four output fields
(\code{general\_search}, \code{news\_search}, \code{person\_linkedin\_scrape},
\code{company\_linkedin\_scrape}) are read by name inside the prompts of every email
workflow.

So the handoff is: run research by hand, let it write to the sheet, then run the
orchestrator by hand, which reads the sheet. Two manual steps and an undocumented
ordering constraint, with nothing in the system enforcing that research ran first. If
someone runs the orchestrator on a fresh lead, the prompts get empty research fields and
the LLM writes generic copy, with no error anywhere.

\textbf{This handoff is inferred}, from the matching field names, not confirmed by any
node connection. There is no wire. The project's own README states this inference
explicitly and honestly, and this reading independently confirms it.
\end{honest}

\begin{honest}{No deduplication of leads, anywhere}
Nothing in the ten files checks whether a lead was already researched, already written
to, or already pushed to Instantly. Running the orchestrator twice on the same list
generates and uploads the sequence twice. For a cold email system, whose entire premise
is not annoying strangers, that is a consequential gap.
\end{honest}

\section{Cross-cutting patterns}
\label{sec:patterns}

Having read all 80 files, these are the habits that recur. They are more useful than any
single walkthrough, because they are what the author would actually be asked to defend.

\subsection{Google Sheets as the database}

124 Google Sheets nodes and 5 Sheets triggers. Sheets is used here as: the lead
database, the job queue, the state store, the log, and the human review interface.

\begin{keyidea}{This is smarter than it first looks}
The instinct is to call this amateurish. Resist it for a moment, because a spreadsheet
buys four things a real database does not:

\begin{itemize}
  \item \textbf{The client can see it.} No admin panel to build.
  \item \textbf{The client can edit it.} Human-in-the-loop review, for free.
  \item \textbf{Zero provisioning.} No migrations, no hosting, no connection strings.
  \item \textbf{It is already the incumbent.} The work was in a spreadsheet before the
    automation existed. Meeting the business where it is beats a migration.
\end{itemize}

For a small agency automating its own chores, this is a defensible engineering decision,
not a shortcut. Saying so, and then saying where it breaks, is a much stronger answer
than either defending it blindly or apologizing for it.
\end{keyidea}

\begin{honest}{Where Sheets-as-database breaks, concretely}
\begin{itemize}
  \item \textbf{No transactions.} A workflow that reads a row, decides, and writes back
    has a race window. Two concurrent runs can both read "not yet emailed" and both
    email.
  \item \textbf{No unique constraint.} This is \emph{why} the dedup problems in
    Section~\ref{sec:coldemail} exist. A database would have refused the duplicate.
  \item \textbf{Rate limits.} The Sheets API caps requests per minute per user. A batch
    job hits it.
  \item \textbf{It does not scale.} \file{cluster-seo-post-generator} reads the entire
    post history on every run to do internal linking. Fine at 50 rows. At 5{,}000, that
    is a very large prompt and a very slow run.
\end{itemize}
The honest position: right call at this size, with a known and stated ceiling.
\end{honest}

\subsection{Polling for slow jobs}

33 \code{wait} nodes exist because image and video generation take minutes, and the APIs
are asynchronous: you submit a job, get a ticket, and must ask later whether it is done.

\begin{figure}[H]
\centering
\begin{tikzpicture}[node distance=8mm]
  \node[box, text width=1.9cm] (sub) {POST job};
  \node[box, right=of sub, text width=1.9cm] (w) {Wait\\\scriptsize n seconds};
  \node[box, right=of w, text width=1.9cm] (chk) {GET status};
  \node[dec, right=9mm of chk] (d) {done?};
  \node[box, right=9mm of d, text width=1.9cm] (use) {use result};
  \draw[flow] (sub) -- (w);
  \draw[flow] (w) -- (chk);
  \draw[flow] (chk) -- (d);
  \draw[flow] (d) -- node[lbl,above]{yes} (use);
  \draw[flow] (d.south) |- ++(0,-8mm) -| node[lbl,below,pos=0.25]{no, wait again} (w.south);
\end{tikzpicture}
\caption{The polling loop, used across the video family.}
\label{fig:poll}
\end{figure}

\begin{honest}{None of these loops has a maximum attempt count}
The polling loops here check "is it done?" and, if not, wait and check again. What none
of them do is count. If a render job fails in a way that leaves it permanently
not-done, the loop spins until n8n's own execution timeout kills it, burning API calls
the whole way. A bounded loop with a failure branch is the standard shape, and it is
absent.
\end{honest}

\subsection{Trusting the LLM's output shape}

The single most common bug class in this repository. The pattern is always the same: a
prompt says "return only JSON" or "return only the quote", and nothing downstream checks
that it did.

n8n \emph{has} a structured output parser node, which forces the model's response into a
declared schema and fails loudly when it does not fit. It appears \textbf{3 times in
1{,}645 nodes}, in just 2 of the 80 files (twice in
\file{deep-research-report-generator}, once in \file{youtube-agent-starter-template}).
In the other 78 files, output shape is a matter of hoping.

\begin{honest}{The verified worst case}
\file{competitor-content-gap-analyzer}'s \code{Clean Content} node has this as its
stored JavaScript source. Read the letters carefully:

\begin{lstlisting}[language=]
if ($input.item.json.body){nnnn$input.item.json.content =
  $input.item.json.body.replaceAll('/^s+|s+$/g', '')...
\end{lstlisting}

The newlines that should separate the statements have been collapsed into literal letter
\code{n} characters, glued onto the next token. So the parser sees a variable called
\code{nnnn\$input}, which does not exist, and throws a \code{ReferenceError} the instant
the \code{if} branch is taken. Separately, the two regular expressions are passed as
\emph{quoted strings} (\code{'/\^{}s+|s+\$/g'}) rather than as regex literals, so even
with the newlines restored they would match nothing.

The node is dead. It has never cleaned anything.

And here is the part that matters: the node is configured with
\code{continueOnFail: true} and \code{alwaysOutputData: true}. So the error is swallowed
whole. The item passes through unmodified, \code{content} and \code{contentShort} are
never set, the downstream InfraNodus call receives effectively nothing, and \emph{the
workflow reports success}. Every run of this automation has produced a content gap
analysis based on no content.

The lesson is not "someone mangled a code node". It is that
\code{continueOnFail: true} turned a loud crash into a silent wrong answer.
Error-tolerance flags are set on 27 nodes in this repository, and each one is a place
where a failure becomes invisible.
\end{honest}

\subsection{Sticky notes as documentation}

186 sticky notes, more than any other node type. Someone genuinely documented these
canvases as they built them.

\begin{honest}{Documentation drifts, and here it demonstrably has}
In \file{calcom-trigger-workflow}, the sticky notes label two reminder stages
"5-Minute Alert" and "Call Started Notification". The actual \code{If} conditions and
the email copy in those same nodes put one at roughly one hour before the call and the
other in the final ten minutes.

A reader who trusts the sticky notes gets the timing wrong. Canvas comments rot exactly
like code comments, and there is nothing to check them.
\end{honest}

\subsection{Orphaned nodes: the visual tool's characteristic debt}

An orphan is a node with no path from any trigger. It cannot ever run. In text code an
unreachable function is at least invisible; on a canvas, an orphan sits there looking
exactly like a working step.

Twelve of the 80 files contain orphans. Verified by computing reachability from every
trigger:

\begin{center}
\begin{tabular}{@{}lrp{6.3cm}@{}}
\toprule
\textbf{File} & \textbf{Orphans} & \textbf{What is stranded} \\
\midrule
\file{whatsapp-ai-sales-lead-system}      & 13 & An entire agent, model, memory and vector store \\
\file{cold-email...Generate\_All\_Emails} & 9  & Old copies of the sub-workflow calls \\
\file{social-media-trend-monitor}         & 8  & Four scraper agents and their merge \\
\file{vapi-call-trigger-utility}          & 8  & An agent, a model, a loop \\
\file{ai-proposal-generator-system}       & 4  & The whole Google Slides output path \\
\file{faceless-shorts-system}             & 4  & An alternative video provider branch \\
\file{booking-time-utility-set} (x3)      & 3--4 & A Calendly path, in all three files \\
\file{presentation-generator-tool}        & 2  & Two of the three slide generators \\
\file{social-media-poster}                & 2  & The YouTube Shorts path \\
\file{fb-ad-creative-generator}           & 1  & A folder-creation step \\
\bottomrule
\end{tabular}
\end{center}

\begin{honest}{The proposal generator is the clearest case}
\file{ai-proposal-generator-system} has two form triggers. One is wired to the PandaDoc
path and works. The other, feeding the Google Slides path, has this in the connections
map:

\begin{lstlisting}[language=]
"On form submission": { "main": [ [ ] ] }
\end{lstlisting}

An empty target list. The key exists, the wire does not. Four nodes (Replace Text,
OpenAI, Google Drive, Gmail) are unreachable from any trigger.

That is worth dwelling on, because it is a trap for anyone auditing these files: a naive
check of "is this trigger in the connections map?" answers \emph{yes} and is
\emph{wrong}. You have to look at whether the list is empty. This document's reachability
analysis was corrected after making exactly that mistake.

And in that stranded path sits a node that always writes the literal cost \code{\$1,850}
into the deck, regardless of what the form said, plus two Gmail nodes whose subject and
body are hardcoded for one specific past client ("Re: Proposal for LeftClick", "Hey
Nick"). If that trigger were ever reconnected without reading the nodes, it would send
every client the same price and the same name.
\end{honest}

\section{Security posture}
\label{sec:security}

An n8n export is a notorious place to leak a secret. This repository was audited for
that specifically, and the result is good enough to state precisely.

\subsection{How n8n stores credentials, and why that saves this project}

A node does not contain its API key. It contains a \emph{pointer} to an entry in n8n's
encrypted credential store:

\begin{lstlisting}[language=,caption={The credential reference on a node. This is the
whole of it. There is no field here that could hold a secret.}]
"credentials": {
  "stripeApi": {
    "id":   "hR3xK9...",     <-- which entry in the credential store
    "name": "Stripe account"  <-- its display label
  }
}
\end{lstlisting}

Verified across all 80 files: there are \textbf{373 credential references}, spanning
\textbf{27 distinct credential types}, and every single one contains exactly two keys,
\code{id} and \code{name}. Not one carries a value.

\begin{keyidea}{Why this matters more than it sounds}
This is a property of n8n's design, not of the author's diligence, and it is the reason
exporting workflows is safe at all. The secrets never enter the JSON, so they cannot
leak through it.

What the author \emph{is} responsible for is everything n8n does not protect: keys typed
into an HTTP node's header field by hand, endpoints, chat IDs, spreadsheet IDs, email
addresses. That is where an export normally leaks, and that is what the audit below
checked.
\end{keyidea}

\subsection{The audit, and its result}

\begin{center}
\begin{tabular}{@{}lp{9.0cm}@{}}
\toprule
\textbf{Checked for} & \textbf{Result} \\
\midrule
Provider key formats & Scanned all 268 tracked files for OpenAI, Google, Slack, GitHub, Telegram, Apify and JWT key shapes. \textbf{Three hits, all placeholders} (\code{Bearer YOUR\_API\_KEY\_HERE} and similar). No live key. \\
Credential values    & 373 references, all \code{id}/\code{name} only. \textbf{Clean.} \\
Bearer tokens        & Six distinct literals, every one a placeholder. \textbf{Clean.} \\
Private endpoints    & Every host referenced is a public vendor API (Google, Stripe, Tavily, Apify, Runware, Novita, Cal.com, ...). No Supabase project URL, no n8n instance hostname. \textbf{Clean.} \\
Email addresses      & All placeholders (\code{@example.com}) except the owner's own business address. \textbf{Clean.} \\
Owner identity       & The \code{shared} block, which normally carries the n8n account name and email, is scrubbed in all 30 files that have one. \textbf{Clean.} \\
\bottomrule
\end{tabular}
\end{center}

\textbf{Conclusion: no live credential is committed to this repository.} That is a real
result, and it was not free: n8n protects the credential store, but the author had to
find and scrub the hardcoded keys, the account names and the emails by hand. The project
README states that live keys were found and stripped, and this audit finds no survivor.

\begin{honest}{The scrub was done by hand, and it shows}
The same field, the n8n account name inside the \code{shared} block, has been replaced
with \textbf{sixteen different placeholder spellings} across the files:
\code{YOUR\_EMAIL\_HERE}, \code{owner@example.com}, \code{Portfolio Owner},
\code{It Expert}, \code{YOUR\_N8N\_PROJECT\_NAME}, \code{REDACTED\_USER},
\code{Project Owner}, and ten more variants.

The security outcome is correct: nothing leaked. But sixteen spellings of one
substitution is the signature of a manual pass, and a manual pass is one that can miss a
file. A scripted scrub with a single canonical placeholder would be both tidier and
genuinely safer, because it could be re-run over the repository to prove nothing was
missed. As it stands, the proof is this audit, which is a snapshot.
\end{honest}

\begin{honest}{Webhook paths are in the files, and that is a judgment call, not a bug}
Each webhook node carries its \code{path}, a UUID. Combined with the n8n instance's
hostname, that is the live URL of the endpoint. The hostname is not in these files, so
the UUID alone opens nothing, and several of these webhooks have no authentication on
them (\file{birthday-greetings-automation}'s README says so plainly).

The risk is therefore low but not zero, and it rests entirely on the instance hostname
staying private. Regenerating the paths on import would cost nothing and remove the
question. This document does not reproduce any of them.
\end{honest}

\section{The honest summary}
\label{sec:honest}

Everything in this section was verified against the exports. Nothing here is taken from
the READMEs on trust, although the READMEs deserve real credit: they found most of it
first.

\subsection{What is genuinely good}

\begin{itemize}
  \item \textbf{The documentation discipline is exceptional.} 59 folders, 59 READMEs, 59
    credential manifests, nine consistent headings, 9{,}304 lines of prose. Plus 186
    sticky notes inside the canvases. This is far above the norm for automation work.
  \item \textbf{The self-criticism is real, and it is specific.} The \emph{Challenges}
    sections are not marketing. They say "the content-cleaning code is broken", "the
    workflow always writes \$1,850", "custom fields are read by array index". Spot-checks
    against the JSON confirmed these claims exactly, including the subtle ones. A README
    that tells you where the body is buried is worth ten that do not.
  \item \textbf{No credentials leaked.} Section~\ref{sec:security}.
  \item \textbf{The sub-workflow architecture, where it is used, is correct.} The
    personal assistant and the Upwork pipeline are properly factored.
  \item \textbf{The hard parts were not dodged.} 128 hand-written API calls, an async
    polling pattern, a RAG pipeline with a considered chunking strategy and a
    delete-before-insert. This is not drag-and-drop work.
\end{itemize}

\subsection{What is weak, ordered by what it would actually cost}

\begin{enumerate}
  \item \textbf{No idempotency on the invoicing workflow.} A retried webhook
    double-bills a real customer. Highest severity in the repository, and the only one
    that costs money on the first occurrence.
  \item \textbf{Silent failures via \code{continueOnFail}.} 27 nodes. The
    \file{competitor-content-gap-analyzer} case proves the failure mode is not
    theoretical: a broken node produced wrong answers and reported success, indefinitely.
  \item \textbf{Copy-paste architecture.} The duplicated RAG ingestion branch, the six
    duplicated copywriting blocks, the nine orphaned call nodes. This is the
    characteristic debt of visual tools and it is present throughout.
  \item \textbf{Orphaned nodes in 12 of 80 files.} 55-plus nodes that cannot run,
    indistinguishable on the canvas from ones that can.
  \item \textbf{Index-based field access.} \code{custom\_fields[19]} is a silent
    time bomb tied to someone else's column order.
  \item \textbf{Unbounded polling loops.} No attempt cap anywhere in the video family.
  \item \textbf{Hardcoded values where expressions belong.} \code{\$1,850}, "1st
    Deposit", "Hey Nick", a Telegram chat ID, the literal caption \code{"hello"} in
    \file{auto-video-generator}.
  \item \textbf{No output validation on LLM responses.} The structured output parser
    node appears 3 times in 1{,}645 nodes, across just 2 of the 80 files.
  \item \textbf{Documentation drift.} Sticky notes disagreeing with the nodes they label.
\end{enumerate}

\subsection{The maturity question, stated plainly}

\begin{honest}{The one thing to get straight before any interview}
28 manual triggers. 9 of 80 workflows active. 12 files with orphaned nodes. 45 disabled
nodes.

This is a portfolio of \textbf{59 designed and built automations}, a meaningful fraction
of which are at demonstration maturity rather than unattended-production maturity. The
repository does not contain the evidence needed to claim otherwise, because execution
logs are not in it, and the \code{active} flags say what they say.

That is still a strong claim. "I designed and built 59 automations across seven problem
domains, documented every one including its faults, and can tell you exactly which are
production-hardened and which are not" is a better answer than "59 automations run my
client's business", and it has the advantage of being defensible against someone who
opens the files.

The failure mode to avoid is being caught claiming the stronger version.
\end{honest}

\subsection{What was inferred rather than confirmed}

In keeping with the rule that inference must be labelled:

\begin{itemize}
  \item \textbf{The taxonomy in Figure~\ref{fig:taxonomy} is this document's own
    invention.} The repository does not group its folders. The seven families are a
    reading of what the automations do; the counts of folders are exact, the boundaries
    are a judgment call.
  \item \textbf{The research-to-orchestrator handoff in the cold email system is
    inferred} from matching field names. There is no connection in the JSON. The
    project's own README says the same and is equally explicit about it.
  \item \textbf{"Built and demonstrated rather than deployed" is an inference} from
    trigger types and active flags. It is a strong inference, but the alternative
    explanation, that workflows were routinely deactivated before export, cannot be
    ruled out from the files alone. The owner knows which is true; this document cannot.
  \item \textbf{Whether the disabled \code{File Updated} trigger in the RAG pipeline was
    switched off deliberately or abandoned mid-refactor cannot be determined} from the
    export. Both readings fit.
  \item \textbf{Whether the assistants actually work well is untestable here.} The
    graphs are sound; the behaviour lives in prompts, and prompts cannot be evaluated by
    reading them.
\end{itemize}

\section{What an interviewer will ask}

Five questions this repository invites, and the honest answer to each.

\textbf{"Is this really engineering, or did you drag boxes?"} 128 hand-written API calls
where no node existed. An async polling pattern. A RAG pipeline that summarizes
spreadsheets before embedding them, because embedding rows produces meaningless chunks.
A sub-workflow architecture with a transcription seam. The boxes cover the easy 80\%; the
interesting part is the rest. And the failures in Section~\ref{sec:honest} are ordinary
software failures, not tool failures, which is itself the argument.

\textbf{"Why is a spreadsheet your database?"} Because the client can see it and edit it,
it needed no provisioning, and the work was already there. It breaks on concurrency,
uniqueness and scale, and the duplicate-lead problem in the cold email system is exactly
that bill coming due. Right call at that size, known ceiling.

\textbf{"Your own README says a node is broken. Why is it still broken?"} The honest
answer is that finding it and fixing it are different jobs, and the repository's purpose
was to document what exists rather than to repair it. The stronger version of the answer
is that \code{continueOnFail} is what made it survivable long enough to become a
documentation problem instead of an outage.

\textbf{"Which of these actually run?"} Nine are marked active. Say the number, then say
which ones, then say why the rest are not. Do not let this one be discovered.

\textbf{"What would you change first?"} Idempotency on the invoice workflow, because it
is the only one that costs money on the first mistake. Then delete
\code{continueOnFail} from anything whose silent failure produces a wrong answer rather
than a skipped item. Then factor the duplicated RAG branch, because it is the change
that gets more expensive every week it is not made.

\vfill
\begin{center}
\rule{0.4\textwidth}{0.4pt}\\[0.3em]
{\small All counts in this document were produced by parsing the 80 exported workflow
files directly. Where the project's own documentation and the exports disagreed, the
exports were treated as authoritative; no such disagreement was found on any claim
spot-checked.}
\end{center}

\end{document}
